% !TEX program = lualatex
\documentclass[12pt, a4paper]{article}


\usepackage{fontspec}
\setmainfont{CMU Serif}
\usepackage[english, russian, ukrainian]{babel}

\usepackage{amsmath, amssymb}
\usepackage{geometry}
\geometry{margin=2.5cm}
\usepackage{hyperref}
\hypersetup{colorlinks=true, linkcolor=blue, urlcolor=blue}
\usepackage{microtype}
\usepackage{parskip}
\usepackage{booktabs}
\usepackage{enumitem}
\usepackage{xcolor}
\usepackage{mdframed}
\usepackage{tabularray}
\UseTblrLibrary{booktabs}
\definecolor{boxbg}{RGB}{245,245,250}
\definecolor{accent}{RGB}{30,80,160}
\definecolor{planckblue}{RGB}{31,119,180}
\definecolor{theoryred}{RGB}{214,39,40}

\usepackage{pgfplots}
\pgfplotsset{compat=1.18}
\usepackage{pgfplotstable}



\newmdenv[backgroundcolor=boxbg, linecolor=accent, linewidth=1pt,
          innerleftmargin=10pt, innerrightmargin=10pt,
          innertopmargin=8pt, innerbottommargin=8pt]{keybox}


\title{\textbf{Конспект: від інфлатона до CMB}\\
       \large Підсумок розмови про основи інфляційної космології}
\author{}
\date{}

% ---------------------------------------------------------------
% Planck 2018 binned TT data
% Columns: ell  Dl[muK^2]  err_minus  err_plus
% Source: Planck Collaboration 2018, arXiv:1807.06209
% ---------------------------------------------------------------
\pgfplotstableread{plankdata.txt}\planckdata

% ---------------------------------------------------------------
% LCDM theory curve (CAMB, Planck 2018 best-fit parameters)
% H0=67.36, Ob h^2=0.02237, Oc h^2=0.1200,
% tau=0.0544, As=2.1e-9, ns=0.9649
% ---------------------------------------------------------------
\pgfplotstableread{LCDM.txt}\theorydata

\begin{document}
\maketitle
\tableofcontents
\newpage

% -------------------------------------------------------
\section{Інфлатон і потенціал}

Інфлатон~--- гіпотетичне скалярне поле $\phi$, яке існувало на початку Всесвіту.
У кожній точці простору є одне число~--- значення поля.
Динаміка визначається потенціалом $V(\phi)$: поле «котиться» з гірки вниз,
як м'яч на схилі.

\begin{keybox}
\textbf{Рівняння руху інфлатона:}
\[
  \ddot{\phi} + 3H\dot{\phi} + V'(\phi) = 0
\]
\textbf{Енергія інфлатона:}
\[
  \rho_\phi = \frac{\dot{\phi}^2}{2} + V(\phi)
\]
Під час інфляції $\dot{\phi}^2 \ll V(\phi)$, тому $\rho_\phi \approx V(\phi) = \text{const}$.
\end{keybox}

Ця майже стала енергія через рівняння Фрідмана дає
\[
  H^2 = \frac{8\pi G}{3}\rho_\phi \approx \text{const}
  \quad\Rightarrow\quad
  a(t) \sim e^{Ht},
\]
тобто \textbf{експоненційне розширення}.

\subsection{Форма потенціалу}

Точна форма $V(\phi)$ невідома. CMB дозволяє відновити лише невеликий відрізок
потенціалу~--- той що відповідає останнім $\sim\!7$--$8$ e-folds інфляції.
Існує понад 100 конкуруючих моделей (Стaробінський, хіггсова інфляція,
натуральна інфляція, хаотична інфляція Лінде тощо).

% -------------------------------------------------------
\section{Квантові флуктуації і збурення}

Поверх класичного фону є квантові флуктуації двох типів:

\begin{enumerate}
  \item \textbf{Флуктуації інфлатона} $\delta\phi$~--- скалярні збурення.
  \item \textbf{Флуктуації метрики} $h_{ij}$~--- тензорні збурення (первинні гравітаційні хвилі).
\end{enumerate}

Амплітуда замороженої флуктуації:
\[
  \delta\phi \sim \frac{H}{\dot{\phi}}
\]
Коли $\dot{\phi} \to 0$ (м'яч на «полиці» потенціалу)~--- амплітуда зростає.
Саме звідси ідея \textbf{зламаного спектра}: сходинка в потенціалі дає
локальне підсилення потужності на відповідному масштабі $k_\mathrm{break}$.

\subsection{Заморожування і розморожування}

Флуктуація заморожується, коли її довжина хвилі перевищує горизонт Хаббла
$\lambda > H^{-1}$. Після інфляції горизонт росте швидше за масштаби~---
флуктуації \textbf{повертаються в горизонт} і починають еволюціонувати.

% -------------------------------------------------------
\section{Метрика Фрідмана і збурення}

Фонова метрика FLRW у конформному часі $\eta$ ($d\eta = dt/a$):
\[
  ds^2 = a(\eta)^2\bigl(-d\eta^2 + dx^2 + dy^2 + dz^2\bigr)
\]
де $a(\eta)$~--- масштабний фактор, динаміку якого задають рівняння Фрідмана:
\[
  \mathcal{H}^2 \equiv \left(\frac{a'}{a}\right)^2 = \frac{8\pi G}{3}\,a^2\bar{\rho},
  \qquad
  \mathcal{H}' = -\frac{4\pi G}{3}\,a^2(\bar{\rho} + 3\bar{p}),
\]
де штрих означає похідну по $\eta$, $\mathcal{H} = aH$~---
конформний параметр Хаббла, $\bar{\rho}(\eta)$~--- однорідна фонова густина, $\delta\rho$~--- збурення:
\[
  \rho(\eta, \mathbf{x}) = \bar{\rho}(\eta) + \delta\rho(\eta, \mathbf{x}),
  \qquad |\delta\rho| \ll \bar{\rho}.
\]
Домінуюча компонента визначає закон розширення:

\begin{center}
\begin{tabular}{llll}
\toprule
Епоха & Джерело & $w = p/\rho$ & $a(\eta)$ \\
\midrule
Інфляція     & Інфлатон                & $\approx -1$ & $\propto e^{Ht}$ \\
Радіаційна   & Фотони, нейтрино        & $1/3$        & $\propto \eta$ \\
Матеріальна  & Темна матерія + баріони & $0$          & $\propto \eta^2$ \\
Сьогодні     & Темна енергія           & $\approx -1$ & $\propto e^{Ht}$ \\
\bottomrule
\end{tabular}
\end{center}

\subsection*{SVT-розклад збурень метрики}

Збурення метрики $\delta g_{\mu\nu}$ розкладаються на три
\emph{незалежних} типи за поведінкою під просторовими обертаннями
(теорема Стюарта--Вокера):
\[
  \delta g_{\mu\nu} =
    \underbrace{\delta g_{\mu\nu}^{(S)}}_{\text{скалярна}: \Phi,\,\Psi}
  + \underbrace{\delta g_{\mu\nu}^{(V)}}_{\text{векторна: затухає}}
  + \underbrace{\delta g_{\mu\nu}^{(T)}}_{
      \text{тензорна: } h_{ij},\; h_{ij}'' + 2\mathcal{H}h_{ij}' - \nabla^2 h_{ij} = 0
    }
\]
Повна збурена метрика в конформно-ньютонівській калібровці:
\[
  ds^2 = a^2\bigl[
    \underbrace{-(1+2\Phi)}_{\text{час: гравіт. потенціал}}\,d\eta^2
    + \bigl(
        \underbrace{(1-2\Psi)\,\delta_{ij}}_{\text{ізотропне стиснення}}
      + \underbrace{h_{ij}}_{\text{гравіт. хвилі}}
      \bigr)dx^i dx^j
  \bigr].
\]

Три частини метрики мають різну фізику і еволюціонують \emph{незалежно}
в лінійній теорії:
\begin{itemize}
  \item $\Phi$~--- ньютонівський гравітаційний потенціал,
        уповільнює (або прискорює) фотони при вході/виході з ям~---
        ефект Закса--Вольфа.
  \item $\Psi$~--- кривина просторових поверхонь,
        визначає геодезичне відхилення частинок.
        За відсутності анізотропного напруження $\Phi = \Psi$
        (виконується до рекомбінації з хорошою точністю).
  \item $h_{ij}$~--- тензорне збурення, поперечне і безслідове
        ($\partial^i h_{ij} = 0$, $h^i{}_i = 0$),
        два поляризаційних стани $h_+$ і $h_\times$.
        Задовольняє рівняння вільної хвилі:
        \[
          h_{ij}'' + 2\mathcal{H}h_{ij}' - \nabla^2 h_{ij} = 0.
        \]
        Це первинні гравітаційні хвилі~--- народжуються інфляцією
        незалежно від скалярних збурень і залишають слід у
        B-модах поляризації CMB.
\end{itemize}

\subsection*{Первинний спектр із інфляції}

Інфляція генерує майже гаусові, майже інваріантні за масштабом
скалярні збурення кривини $\mathcal{R}_k$ з амплітудою:
\[
  \mathcal{P}_{\mathcal{R}}(k)
    \equiv \frac{k^3}{2\pi^2}\,|\mathcal{R}_k|^2
    = A_s \left(\frac{k}{k_*}\right)^{n_s-1},
\]
де $k_* = 0.05\,\mathrm{Мпк}^{-1}$~--- опорний масштаб,
$A_s$~--- амплітуда, $n_s < 1$~--- спектральний індекс
(відхилення від Harrison--Zeldovich спектра).
Саме цей спектр є початковою умовою для всієї подальшої еволюції збурень.


% -------------------------------------------------------
\section{Реогрів і народження матерії}

Коли інфлатон докочується до мінімуму потенціалу, він починає коливатись
і розпадається на частинки~--- відбувається \textbf{реогрів}.
Де енергія інфлатона була більшою~--- там народилось більше частинок.
Так флуктуації $\delta\phi$ передаються матерії як \textbf{початкові умови}
для нерівномірного розподілу густини.

\begin{keybox}
Після реогріву інфлатона більше немає.
Гравітаційні хвилі $h_{ij}$ продовжують існувати незалежно~---
вони майже не взаємодіють з матерією.
\end{keybox}

% -------------------------------------------------------
\section{Акустичні коливання і CMB}

Після реогріву~--- гаряча плазма фотонів і барионів.
Нерівномірний розподіл густини від інфлатона стає початковим поштовхом.
Тиск фотонів і гравітація породжують \textbf{акустичні коливання}~---
стоячі хвилі в плазмі.

При рекомбінації ($z \approx 1100$, вік $\approx 380\,000$ років)
плазма охолола, фотони полетіли~--- коливання «заморозились».
Масштаби що встигли зробити ціле число півперіодів дали \textbf{піки} на CMB.

\subsection{Спектр потужності CMB}

Зв'язок між первинним спектром і спостережуваним:
\[
  C_\ell = \int_0^\infty \frac{dk}{k}\, P(k)\,\bigl|\Delta_\ell(k)\bigr|^2
\]
де $\Delta_\ell(k)$~--- функція переносу (рахується числово в CAMB/CLASS).

\subsection{Від метрики Фрідмана до кривої CMB}

Метрика FLRW через рівняння Фрідмана визначає два ключових масштаби,
які керують положенням піків.

\textbf{Звуковий горизонт}~--- відстань, яку звукова хвиля пройшла
в плазмі до рекомбінації $\eta_*$:
\[
  r_s(\eta_*) = \int_0^{\eta_*} c_s(\eta)\, d\eta,
  \qquad
  c_s = \frac{c}{\sqrt{3(1 + R)}},\quad
  R = \frac{3\rho_b}{4\rho_\gamma}
\]
де $R$~--- відношення барионної та фотонної густини
(визначається параметром $\Omega_b h^2$).

\textbf{Кутова відстань} до поверхні останнього розсіяння:
\[
  d_A(\eta_*) = \int_{\eta_*}^{\eta_0} \frac{d\eta}{1+z}
\]
залежить від $H_0$, $\Omega_c h^2$, $\Omega_\Lambda$.

\textbf{Положення піків.}
Піки відповідають масштабам, що встигли зробити ціле число
півперіодів коливань до рекомбінації:
\[
  \ell_m \approx m\,\frac{\pi\, d_A(\eta_*)}{r_s(\eta_*)},
  \qquad m = 1, 2, 3, \ldots
\]
Перший пік $\ell_1 \approx 220$ безпосередньо вимірює відношення $d_A/r_s$,
тобто геометрію Всесвіту.

\textbf{Акустичний осцилятор.}
Фотон-барионна рідина в гравітаційній ямі темної матерії описується рівнянням:
\[
  \Theta'' + \frac{\mathcal{H}(1-3c_s^2)}{1+R}\,\Theta'
  + k^2 c_s^2\,\Theta
  = -k^2\!\left(\Psi - \frac{\Phi''}{k^2 c_s^2}\right)
\]
де $\Theta = \delta T/T$~--- температурне збурення,
$\Phi, \Psi$~--- гравітаційні потенціали зі збуреної метрики FLRW.
Звідси якісно:
непарні піки (стиснення в ямі) вищі за парні~--- ефект барионного
завантаження $R \propto \Omega_b h^2$;
темна матерія поглиблює потенціальні ями, підсилюючи всі піки відносно плато.

\textbf{Плато Закса--Вольфа} ($\ell \lesssim 30$).
Фотони при вильоті з потенціальної ями отримують гравітаційне червоне зміщення:
\[
  \frac{\delta T}{T}\bigg|_{\rm SW} \approx \frac{1}{3}\Phi(\eta_*)
  \quad\Rightarrow\quad
  \mathcal{D}_\ell^{\rm SW} \approx
  \frac{A_s}{9}\,\frac{\ell(\ell+1)}{2\pi}
  \left(\frac{k}{k_*}\right)^{n_s-1}\bigg|_{k \sim \ell/\eta_0}
\]
При $n_s = 1$ плато горизонтальне; $n_s < 1$ дає легкий нахил~--- слід інфляції.

\textbf{Загасання Силка} ($\ell \gtrsim 1000$).
Дифузія фотонів розмиває анізотропію на малих масштабах:
\[
  C_\ell \propto e^{-\ell^2/\ell_D^2}, \qquad \ell_D \sim 1500.
\]

\begin{keybox}
\textbf{Два шари інформації в одному знімку CMB:}
\begin{itemize}[nosep]
  \item \textbf{Нахил під піками} ($n_s \approx 0.965$)~--- натяк про інфляцію,
        відбиток форми потенціалу.
  \item \textbf{Самі піки}~--- акустика плазми, пізніша фізика.
\end{itemize}
\end{keybox}

\subsection{Спектр температурних анізотропій CMB}

На рисунку~\ref{fig:cmb_tt} показано кутовий спектр потужності температурних
анізотропій CMB (TT-спектр):
\[
  \mathcal{D}_\ell^{TT} \equiv \frac{\ell(\ell+1)}{2\pi}\,C_\ell^{TT}
  \quad [\mu\mathrm{K}^2]
\]
Точки~--- збінновані дані Planck~2018 з~$1\sigma$ похибками.
Суцільна крива~--- теоретичне передбачення базової моделі $\Lambda$CDM,
розраховане програмою CAMB з параметрами найкращого фіту Planck~2018
(таблиця~\ref{tab:params}).

\begin{figure}[h!]
\centering
\begin{tikzpicture}
\begin{semilogxaxis}[
    width  = 14.5cm,
    height = 8.5cm,
    xlabel = {Мультипольний момент $\ell$},
    ylabel = {$\mathcal{D}_\ell^{TT}\ [\mu\mathrm{K}^2]$},
    xmin   = 2,   xmax = 2600,
    ymin   = 0,   ymax = 6800,
    xtick  = {2, 10, 30, 100, 300, 1000, 2500},
    xticklabels = {2, 10, 30, 100, 300, 1000, 2500},
    minor xtick = {5,20,50,200,500,2000},
    ytick  = {0, 1000, 2000, 3000, 4000, 5000, 6000},
    grid   = both,
    grid style = {line width=0.3pt, draw=gray!25},
    major grid style = {line width=0.5pt, draw=gray!40},
    tick label style = {font=\small},
    label style      = {font=\small},
    legend style = {
        at={(0.6,0.97)}, anchor=north east,
        font=\small, fill=white, fill opacity=0.9,
        draw=gray!60, inner sep=4pt, row sep=1pt
    },
    clip = true,
]

%--- Theory curve (LCDM, CAMB) ---
\addplot[
    color     = theoryred,
    line width = 1.6pt,
    smooth,
] table[x=ell, y=Dl] {\theorydata};
\addlegendentry{$\Lambda$CDM (CAMB, Planck\,2018 best-fit)}

%--- Planck 2018 data with error bars ---
\addplot[
    only marks,
    mark       = *,
    mark size  = 1.4pt,
    color      = planckblue,
    error bars/.cd,
        y dir  = both,
        y explicit,
] table[
    x    = ell,
    y    = Dl,
    y error minus = errm,
    y error plus  = errp,
] {\planckdata};
\addlegendentry{Planck\,2018 (збіновані дані TT)}

%--- Annotate acoustic peaks ---
\node[font=\tiny, text=gray!70, above] at (axis cs:220,5732) {1-й пік};
\node[font=\tiny, text=gray!70, above] at (axis cs:540,2592) {2-й пік};
\node[font=\tiny, text=gray!70, above] at (axis cs:810,2540) {3-й пік};

%--- Sachs-Wolfe plateau annotation ---
\draw[->, gray!60, thin] (axis cs:6,1300) -- (axis cs:30,1070)
    node[midway, left, font=\tiny, text=gray!60] {плато Закса--Вольфа};

\end{semilogxaxis}
\end{tikzpicture}
\caption{Кутовий TT-спектр потужності CMB. Точки~--- дані Planck~2018
         \cite{planck2018params}.
         Крива~--- базова модель $\Lambda$CDM, розрахована CAMB
         з параметрами таблиці~\ref{tab:params}.
         Видно три акустичні піки та плато Закса--Вольфа на великих
         кутових масштабах ($\ell \lesssim 30$).}
\label{fig:cmb_tt}
\end{figure}

\begin{table}[h!]
\centering
\caption{Параметри базової моделі $\Lambda$CDM~--- найкращий фіт Planck~2018
         TT,TE,EE+lowE+lensing \cite{planck2018params}.}
\label{tab:params}
\begin{tblr}{
colspec={lllX},
row{1}={c,m,font=\bfseries}
}
\toprule
Параметр & Позначення & Значення & Фізичний зміст \\
\midrule
Стала Хаббла           & $H_0$              & $67.36 \pm 0.54$ км/с/Мпк
                        & темп розширення сьогодні \\
Барионна густина       & $\Omega_b h^2$     & $0.02237 \pm 0.00015$
                        & частка барионної матерії \\
Темна матерія          & $\Omega_c h^2$     & $0.1200  \pm 0.0012$
                        & частка холодної темної матерії \\
Оптична глибина        & $\tau$             & $0.0544  \pm 0.0073$
                        & реіонізація Всесвіту \\
Амплітуда спектра      & $\ln(10^{10}A_s)$  & $3.044   \pm 0.014$
                        & «гучність» флуктуацій \\
Спектральний індекс    & $n_s$              & $0.9649  \pm 0.0042$
                        & нахил інфляційного спектра \\
\bottomrule
\end{tblr}
\end{table}

\begin{keybox}
\textbf{Як параметри $\Lambda$CDM узгоджуються з кривою:}
\begin{itemize}
  \item $n_s < 1$~--- нахил спектра на великих масштабах (плато не горизонтальне).
  \item $\Omega_b h^2$~--- визначає відносну висоту парних і непарних піків
        (баріони підсилюють стиснення, гальмують розрідження).
  \item $\Omega_c h^2$~--- зсуває положення піків і визначає форму між ними
        (темна матерія прискорює початок матеріальної епохи).
  \item $H_0$~--- визначає кутовий масштаб звукового горизонту,
        тобто положення першого піку по~$\ell$.
  \item $\tau$~--- гасить спектр при $\ell \gtrsim 10$ на рівні $e^{-2\tau}$.
  \item $A_s$~--- загальна вертикальна нормалізація всієї кривої.
\end{itemize}
\end{keybox}

\begin{thebibliography}{9}
\bibitem{planck2018params}
  Planck Collaboration,
  \textit{Planck 2018 results. VI. Cosmological parameters},
  A\&A \textbf{641}, A6 (2020),
  arXiv:1807.06209.
\end{thebibliography}

%% -------------------------------------------------------
%\section{Параметри інфляції з CMB}
%
%\begin{center}
%\begin{tabular}{lll}
%\toprule
%Параметр & Визначення & Що говорить про потенціал \\
%\midrule
%$A_s \approx 2.1\times10^{-9}$ & Амплітуда спектра & Висота флуктуацій (гучність) \\
%$n_s \approx 0.965$             & Спектральний індекс & Нахил схилу $V(\phi)$ \\
%$r < 0.06$                      & Тензорне відношення & Висота потенціалу \\
%\bottomrule
%\end{tabular}
%\end{center}
%
%Первинний спектр інфляційних збурень:
%\[
%  P(k) = A_s \left(\frac{k}{k_*}\right)^{n_s - 1}
%\]

% -------------------------------------------------------
\section{Поляризація CMB}

Фотони поляризуються при останньому розсіянні на електронах.
Напрямок поляризації залежить від \textbf{поперечного} перепаду температури
навколо точки розсіяння.

Поляризаційне поле розкладають на два типи:

\begin{center}
\begin{tabular}{lll}
\toprule
Тип & Джерело & Статус \\
\midrule
E-моди (симетричний візерунок) & Скалярні флуктуації & Виміряні \\
B-моди (закручений візерунок)  & Тензорні флуктуації (GW) & Поки не знайдені \\
\bottomrule
\end{tabular}
\end{center}

Знайти B-моди~--- означає отримати прямий доказ первинних гравітаційних хвиль
і квантової природи гравітації в режимі малих збурень.

% -------------------------------------------------------
\section{Хаотична інфляція Лінде}

Лінде (1983): початкове значення $\phi$ хаотичне~--- різне в різних точках простору.
Найпростіший потенціал: $V(\phi) = \frac{1}{2}m^2\phi^2$.

Де $\phi$ велике~--- інфляція іде довго, виникає величезна «бульбашка».
Наш Всесвіт~--- одна з таких бульбашок. Інші бульбашки недосяжні~---
простір між ними розширюється швидше за світло.
Це \textbf{вічна інфляція} і мультивсесвіт як математичний наслідок.

% -------------------------------------------------------
\section{Наша задача: зламаний спектр і JWST}

\subsection{Мотивація}

JWST знайшов масивні галактики при $z \sim 10$--$12$, яких у стандартній
моделі бути не повинно~--- не вистачало часу утворитись.

\subsection{Ідея}

Якщо потенціал інфлатона має \textbf{сходинку}, м'яч сповільнюється~---
$\dot{\phi} \to 0$~--- і флуктуації що заморожуються в цей момент мають
більшу амплітуду. Зламаний спектр:
\[
  P(k) = A_s \left(\frac{k}{k_*}\right)^{n_s-1}
  \times \mathcal{F}(k;\, k_\mathrm{break},\, \Delta n)
\]
Більше потужності на малих масштабах $k \gtrsim 1\,\mathrm{Mpc}^{-1}$~---
гало тієї самої маси утворюються \textbf{раніше}~--- узгоджується з JWST.

\subsection{Чому не суперечить CMB і $\sigma_8$}

Planck чутливий до $k \lesssim 0.2\,\mathrm{Mpc}^{-1}$~--- злом не бачить.
$\sigma_8$ визначається масштабом $\sim 8\,\mathrm{Mpc}/h$, тобто
$k \sim 0.1$--$0.5\,\mathrm{Mpc}^{-1}$~--- майже не змінюється.

\subsection{План роботи}

\begin{enumerate}
  \item Огляд літератури: чи зроблено вже саме це поєднання.
  \item Мінімальна модель: 2 нових параметри $(k_\mathrm{break},\, \Delta n)$.
  \item CAMB/CLASS: розрахунок $P(k)$, $C_\ell$, $\sigma_8$.
  \item HMF (\texttt{hmf}, Python): порівняння з каталогами JWST.
  \item MCMC (Cobaya): posterior по повному набору параметрів.
  \item Стаття: чи є сумісна область параметрів для Planck + JWST + $\sigma_8$.
\end{enumerate}

\end{document}
